% !TeX program = xelatex
% ============================================================
% pmi/pmi.tex — Программа и методика испытаний
% ГОСТ 19.301-79 «Программа и методика испытаний.
% Требования к содержанию и оформлению»
% ============================================================
\documentclass[12pt,a4paper]{extarticle}

\input{../common/preamble}

\newcommand{\docname}{Test Program and Methodology}
\newcommand{\doccode}{\hseDocCodeBase\ 51 01-1}
\newcommand{\doccodelu}{\hseDocCodeBase\ 51 01-1-ЛУ}

\input{../common/commands}
\input{../common/styles}

\begin{document}
\sloppy

% ── Титульный лист ──────────────────────────────────────────
\input{../common/title_template}


% ============================================================
%  ENGLISH DOCUMENT BODY (project.lang == en)
%  §2.3.4.2 / §2.5.3.1: a translated GOST 19 document rendered in English.
%  The ГОСТ (ЕСПД) section order is preserved 1:1 with the Russian branch.
%  References follow IEEE style. Fill the yellow placeholders.
% ============================================================

% ============================================================
% ANNOTATION
% ============================================================
\section*{ANNOTATION}
\addcontentsline{toc}{section}{ANNOTATION}

This Test Program and Methodology for the development of the program \enquote{\hseProjectName} contains the following sections: \enquote{Test Object}, \enquote{Test Objective}, \enquote{Requirements for the Program}, \enquote{Requirements for the Program Documentation}, \enquote{Test Facilities and Procedure}, \enquote{Test Methods}.

The section \enquote{Test Object} specifies the name, application area and designation of the program under test.

The section \enquote{Test Objective} specifies the objective of the testing.

The \enquote{Requirements for the Program} section lists the requirements from the statement of work that are subject to verification during testing.

The \enquote{Requirements for the Program Documentation} section lists the program documentation submitted for testing.

The section \enquote{Test Facilities and Procedure} contains information on the hardware and software used during testing, as well as the procedure for carrying it out.

The section \enquote{Test Methods} describes the methods used and the specific checks.

% Hint: list what is provided in the appendices of your Test Program and Methodology.
The appendices provide a list of terms, a traceability matrix of requirements to automated tests and a list of supporting test materials.

\clearpage

% ============================================================
% LIST OF APPLICABLE STANDARDS
% ============================================================
\noindent
This document has been developed in accordance with the requirements of:

\begin{enumerate}[label=\arabic*), leftmargin=1.25cm]
\item GOST 19.101-77 Types of programs and program documents\cite{gost19101}.
\item GOST 19.103-77 Designations of programs and program documents\cite{gost19103}.
\item GOST 19.104-78 Basic inscriptions\cite{gost19104}.
\item GOST 19.105-78 General requirements for program documents\cite{gost19105}.
\item GOST 19.106-78 Requirements for program documents executed
      in printed form\cite{gost19106}.
\item GOST 19.301-79 Test program and methodology. Requirements
      for the content and presentation\cite{gost19301}.
\end{enumerate}

Changes to this document are made in accordance with GOST\,19.604-78\cite{gost19604}.

\clearpage

% ============================================================
% CONTENTS
% ============================================================
\hypersetup{linkcolor=black}
\tableofcontents
\hypersetup{linkcolor=blue}
\thispagestyle{fancy}
\clearpage

% ============================================================
% 1. TEST OBJECT
% ============================================================
\section{TEST OBJECT}

\subsection{Name of the Program Under Test}

The name of the program~--- \enquote{\projectname}.

% The English name is substituted from the project settings (the
% \enquote{Work title in English} field).
The name of the program in English~--- \enquote{\hseProjectNameEn}.

\subsection{Application Area of the Program Under Test}

% Hint: briefly describe the purpose and application area of the program,
% and which system boundary is considered within the scope of the tests.

The program \enquote{\projectname} is \hseFill{a brief description of the system purpose}. The system covers \hseFill{the main domains and subsystems}.

Within this document, \hseFill{the system boundary under test} is considered.

\subsection{Designation of the Program Under Test}

The name of the development topic~--- \enquote{\projectname}.

% The short name is substituted from the project settings (the
% \enquote{Short program name} field).
The short name of the program~--- \enquote{\hseProjectNameShort}.

\clearpage

% ============================================================
% 2. TEST OBJECTIVE
% ============================================================
\section{TEST OBJECTIVE}

The objective of the testing is to verify the conformance of the characteristics of the developed program to the functional and other requirements set out in the program documentation \enquote{Statement of Work}.

\clearpage

% ============================================================
% 3. REQUIREMENTS FOR THE PROGRAM
% ============================================================
\section{REQUIREMENTS FOR THE PROGRAM}

\subsection{Requirements for Functional Characteristics}

\subsubsection{Composition of the Functions Performed}

% Hint: identifiers (ТЗ-Ф-01, etc.) link requirements to tests
% in the \enquote{Requirements traceability} panel.

The functional characteristics of the program listed in the requirements table of the statement of work are subject to verification. Each functional requirement with an identifier of the form \texttt{ТЗ-Ф-XX} is matched to an automated test, a test suite or another check.

\input{../technical_specification/requirements_table}

\clearpage

\subsubsection{Organization of Input Data}

% Hint: list the input data and formats of your system to be verified
% (request structure, validation, constraints).
\begin{hseExample}
During testing, compliance with the statement of work requirements for the organization of input data is verified. The following are subject to verification:

\begin{itemize}[leftmargin=1.25cm]
    \item acceptance of \hseFill{type of requests or input data} with data transfer in the \hseFill{input data format} format;
    \item validation of incoming data against the schemas and data types specified in the interface specification;
    \item compliance with the size limits \hseFill{type of input data limits};
    \item \hseFill{other input data organization checks}.
\end{itemize}

Verification is carried out both on valid input data and on erroneous data sets intended to confirm the operation of validation.
\end{hseExample}

\subsubsection{Organization of Output Data}

% Hint: list the output data to be verified (response codes,
% formats, error structure).
\begin{hseExample}
During testing, compliance with the statement of work requirements for the organization of output data is verified. The following are subject to verification:

\begin{itemize}[leftmargin=1.25cm]
    \item returning results in the \hseFill{output data format} format;
    \item use of \hseFill{response status codes or statuses} for successful and erroneous responses;
    \item returning a structured error description sufficient for diagnostics without disclosing confidential data;
    \item \hseFill{other output data organization checks}.
\end{itemize}
\end{hseExample}

\subsubsection{Requirements for Timing Characteristics}

% Hint: list the timing indicators from the statement of work to be
% verified (response time, throughput).
\begin{hseExample}
During testing, the timing characteristics established in the statement of work are verified:

\begin{itemize}[leftmargin=1.25cm]
    \item the processing time of \hseFill{type of requests or operations} must not exceed \hseFill{target response time};
    \item the specified indicators must be maintained under normal load of up to \hseFill{target number of users};
    \item the target request rate must correspond to the range \hseFill{target request rate}.
\end{itemize}
\end{hseExample}

\subsection{Requirements for the Interface}

% Hint: describe the requirements for the program interface to be verified
% (API methods, contracts, specifications).
\begin{hseExample}
During testing, compliance with the statement of work requirements for the program interface is verified. The following are subject to verification:

\begin{itemize}[leftmargin=1.25cm]
    \item implementation of the \hseFill{type of program interface} in accordance with \hseFill{interface principles and specifications};
    \item availability of the current interface specification \hseFill{location of the specification};
    \item verification of access to protected functions based on \hseFill{authentication and authorization mechanism};
    \item uniformity of response structure, data formats and field naming conventions.
\end{itemize}
\end{hseExample}

\subsection{Requirements for Marking and Versioning}

% Hint: adjust the list to the requirements of your statement of work,
% if they are specified.
\begin{hseExample}
During testing, compliance with the statement of work requirements for the marking and versioning of system components is verified. The following are subject to verification:

\begin{itemize}[leftmargin=1.25cm]
    \item conformance of the program code versions to the \hseFill{versioning scheme used} specification;
    \item presence of version tags \hseFill{type of build artifacts};
    \item presence of build metadata with the build date and the source code version identifier;
    \item storage of source code, configurations and build artifacts in repositories that support version control.
\end{itemize}
\end{hseExample}

\clearpage

% ============================================================
% 4. REQUIREMENTS FOR THE PROGRAM DOCUMENTATION
% ============================================================
\section{REQUIREMENTS FOR THE PROGRAM DOCUMENTATION}

\subsection{Composition of the Program Documentation Submitted for Testing}

% Hint: list the documents submitted for testing, indicating the
% corresponding GOST standards.

The following documentation for the program must be submitted for testing:
\begin{enumerate}[label=\arabic*), leftmargin=1.25cm]
    \item \enquote{\projectname}. Statement of Work (GOST 19.201-78).
    \item \enquote{\projectname}. Source Listing (GOST 19.401-78).
    \item \enquote{\projectname}. Programmer's Manual (GOST 19.504-79).
    \item \enquote{\projectname}. Test Program and Methodology
          (GOST 19.301-79).
\end{enumerate}

\subsection{Special Requirements}

The documentation for the program must be executed in accordance with GOST 19.106-78\cite{gost19106} and the GOST standards for each type of document (see Section~4.1).

% Hint: specify other special formatting requirements for the
% documentation if necessary.
Supporting test materials may be formatted as appendices to this document.

\clearpage

% ============================================================
% 5. TEST FACILITIES AND PROCEDURE
% ============================================================
\section{TEST FACILITIES AND PROCEDURE}

\subsection{Hardware Used During Testing}

% Hint: describe the hardware environment of the tests (OS, processor,
% memory) and the launch environments.
\begin{hseExample}
Testing is carried out on hardware with the following characteristics: \hseFill{OS, processor, memory}. The following test environments are used:

% The table is set with \captionof (not \begin{table}[H]): floating
% environments inside the yellow example block are not allowed.
\begin{center}
\small
\captionof{table}{Test environments}
\label{tab:pmi_test_environments}
\begin{tabularx}{\textwidth}{|>{\raggedright\arraybackslash}p{0.21\textwidth}|>{\raggedright\arraybackslash}p{0.29\textwidth}|X|}
\hline
\textbf{Environment} & \textbf{Composition} & \textbf{Purpose} \\
\hline
Local environment & \hseFill{composition of the local environment} & Quick test runs without fully deploying the environment. \\
\hline
\hseFill{integration environment} & \hseFill{composition of the integration environment} & Full integration checks in an environment as close as possible to the production environment. \\
\hline
\hseFill{acceptance stand} & \hseFill{composition of the acceptance stand} & Acceptance checks of the operational characteristics of the deployed system. \\
\hline
\end{tabularx}
\end{center}
\end{hseExample}

\subsection{Software Used During Testing}

% Hint: list the software and tools (runtime environment, testing
% frameworks, quality control tools).
\begin{hseExample}
The following software is used during testing:

\begin{itemize}[leftmargin=1.25cm]
    \item runtime environment: \hseFill{languages and runtime environments};
    \item testing frameworks: \hseFill{testing frameworks};
    \item static analysis tools: \hseFill{linters and analysers};
    \item quality control and coverage tools: \hseFill{quality control tools};
    \item infrastructure tools: \hseFill{DBMS, brokers, containerization}.
\end{itemize}

The main method of accepting functional requirements in this work is passing automated tests: input data, actions and expected results are fixed in executable test scenarios, and the results of their execution are stored in \hseFill{CI/CD or test reports}.
\end{hseExample}

\subsection{Acceptance Criterion for Functional Requirements}

\begin{hseExample}
A functional requirement with an identifier of the form \texttt{ТЗ-Ф-XX} is considered accepted if all the following conditions are met simultaneously:

\begin{enumerate}[label=\arabic*), leftmargin=1.25cm]
    \item the requirement is present in the requirements table of the statement of work;
    \item an automated test, a test suite or an end-to-end scenario is specified for the requirement;
    \item the corresponding check on the version of the program under test completed successfully;
    \item supporting artifacts are available for the check: an execution log, a test report or a coverage report, if they are provided for by this environment.
\end{enumerate}

The acceptance committee confirms the fulfilment of the functional requirements in one of two ways:

% Hint: refine the confirmation methods for your project.
\begin{enumerate}[label=\arabic*), leftmargin=1.25cm]
    \item request access to the \hseFill{project repository}, open the check results of the version under test and verify their successful completion;
    \item run the tests locally on the agreed source code version using the command \hseFill{test run command}.
\end{enumerate}

A local run is used as an alternative confirmation method when access to CI/CD is unavailable or when independent reproduction of the result is required. If there are discrepancies between the local run and CI/CD, the acceptance result is the run on the agreed version in an isolated CI/CD environment, since it fixes the code version, configuration, environment and execution artifacts.
\end{hseExample}

\subsection{Testing Procedure}

% Hint: adjust the sequence of stages for your project.
\begin{hseExample}
Testing is carried out in the following order:

\begin{enumerate}[label=\arabic*), leftmargin=1.25cm]
    \item \textbf{Preparatory stage}: visual verification of the composition of the program documentation, preparation of the test environment according to Sections 5.1--5.2 and deployment of the program according to the programmer's manual.
    \item \textbf{Functional acceptance}: verification of the statement of work requirements based on the results of the checks according to Section~\ref{sec:methods}.
    \item \textbf{Static quality control}: running and analysing the results of \hseFill{linters and quality checks}.
    \item \textbf{Module- and service-level testing}: running \hseFill{testing commands or tasks} and analysing the test and coverage reports.
    \item \textbf{End-to-end testing}: running end-to-end scenarios in the \hseFill{end-to-end testing environment}.
    \item \textbf{Recording of results}: comparing the obtained results with the expected ones and preparing a conclusion on passing the tests.
\end{enumerate}
\end{hseExample}

\subsection{Quantitative and Qualitative Characteristics Subject to Evaluation}

% Hint: adjust the list of evaluated characteristics for your project
% (observability, load, etc.).
\begin{hseExample}
The following characteristics are evaluated during testing:

\begin{itemize}[leftmargin=1.25cm]
    \item the percentage of successfully executed tests and checks;
    \item the absence of unforeseen failures, crashes and inconsistent data states;
    \item the correctness of \hseFill{key verified aspects of the program};
    \item the presence and quality of the supporting test artifacts;
    \item \hseFill{additional characteristics from the statement of work}.
\end{itemize}
\end{hseExample}

\subsection{Testing Conditions}

% Hint: specify the necessary conditions (access rights, source code
% versions, environment, secrets).
\begin{hseExample}
Testing must be carried out subject to the following conditions:

\begin{itemize}[leftmargin=1.25cm]
    \item access to the \hseFill{project repository} and the check artifacts;
    \item a known identifier of the version under test: a commit identifier, a version tag or a link to a merge request;
    \item the presence of \hseFill{the program runtime environment} and network access to the necessary resources;
    \item prepared environment variables, secrets and keys of the test environments;
    \item running the tests on the agreed source code version and test configuration.
\end{itemize}
\end{hseExample}

\clearpage

% ============================================================
% 6. TEST METHODS
% ============================================================
\section{TEST METHODS}
\label{sec:methods}

\begin{hseExample}
Testing is performed by a combination of visual documentation control, automated test suites, quality control checks and acceptance checks. For each functional requirement, traceability down to a specific test, test suite or check from the test methods table is used.
\end{hseExample}

% Hint: for each requirement, describe the verification method, the expected
% and actual results. Fill in the test methods table.
% The test methods table stays outside the yellow example block:
% a longtable inside hseExample is not allowed.

\begin{longtable}{|p{0.18\textwidth}|p{0.22\textwidth}|p{0.25\textwidth}|p{0.22\textwidth}|}
\caption{Test methods}
\label{tab:test-methods}\\
\hline
\textbf{Requirement under test} &
\textbf{Verification method} &
\textbf{Expected result} &
\textbf{Actual result} \\
\hline
\endfirsthead
\hline
\textbf{Requirement under test} &
\textbf{Verification method} &
\textbf{Expected result} &
\textbf{Actual result} \\
\hline
\endhead
\hline
\endfoot
\hline
\endlastfoot
\reqref{ТЗ-Ф-01} &
\hseFill{verification method} &
\hseFill{expected result} &
\hseFill{actual result} \\
\hline
\end{longtable}

\subsection{Testing of Compliance with the Program Documentation Requirements}

Verification is carried out by the visual method. For each document from Section~4.1, the following are verified:

\begin{itemize}[leftmargin=1.25cm]
    \item the presence of the document in the set;
    \item conformance to the designation, the title page and the section structure;
    \item the absence of references to outdated components, interfaces and test environments;
    \item consistency of this document with the statement of work, the source listing and the programmer's manual.
\end{itemize}

\subsection{Testing of Functional Requirements}

% Hint: adjust the steps for your CI/CD or the way you run the tests.
\begin{hseExample}
Verification is carried out by the automated method according to the requirements-to-tests traceability table. The tester performs the following actions:

\begin{enumerate}[label=\arabic*), leftmargin=1.25cm]
    \item open the \hseFill{project repository or CI/CD};
    \item select the check corresponding to the commit or version tag under test;
    \item make sure that all testing jobs have completed successfully;
    \item for each requirement \texttt{ТЗ-Ф-XX}, match the specified test or test suite with a successfully completed check;
    \item save or attach screenshots and execution reports to the test materials.
\end{enumerate}

The test is considered successful if all functional requirements are traced to an automated check, and the corresponding checks are completed without unresolved test failures. If a particular test was re-run due to environment instability, the final successful run and a link to the re-run log are recorded in the test materials.
\end{hseExample}

\subsection{Testing of Static Quality Control and Security Checks}

% Hint: list the linters, quality analysers and security checks used;
% remove the subsection if they are not applicable.
\begin{hseExample}
Verification is carried out by the automated method. The following groups of checks are used:

\begin{itemize}[leftmargin=1.25cm]
    \item \hseFill{linters and formatters} for analysing formatting, linting and basic quality rules;
    \item \hseFill{quality analysis platform} for building quality analytics and aggregating test coverage;
    \item \hseFill{security checking tools} for finding vulnerabilities in dependencies and secrets in the source code.
\end{itemize}

The result of the test is the presence of the generated reports and the absence of critical blockers that make the acceptance of the version impossible.
\end{hseExample}

\subsection{Testing of Timing Characteristics by the Load Method}

% Hint: this subsection is needed if the statement of work specifies timing
% characteristics; otherwise remove it.
\begin{hseExample}
The verification confirms the timing characteristics declared in the statement of work. The testing tool is \hseFill{load testing tool}.

The order of the test:
\begin{enumerate}[label=\arabic*), leftmargin=1.25cm]
    \item prepare the workstation and deploy the program under test in the test environment;
    \item perform a load run using the command \hseFill{run launch command};
    \item perform a stepwise series of runs with the parameters \hseFill{number of users and duration}; consider the first \hseFill{warm-up period duration} as the warm-up period and exclude it from the acceptance metrics;
    \item record the summary indicators of the run: the total number of requests, the number of failures, the aggregated request rate, as well as the response time percentiles (p50, p95, p99);
    \item save or attach the report and the screenshots of the summary statistics to the test materials.
\end{enumerate}

Pass criteria:
\begin{itemize}[nosep, leftmargin=1.25cm]
    \item the aggregated request rate after the warm-up period is at least \hseFill{target request rate};
    \item the proportion of successful responses after the warm-up period is not lower than \hseFill{target proportion of successful responses};
    \item there are no unhandled exceptions and service crashes during the run;
    \item a report is generated with a request table, response time statistics and load graphs.
\end{itemize}
\end{hseExample}

\subsection{Testing of Requirements Traceability and Completeness of Artifacts}

\begin{hseExample}
For each functional requirement of the statement of work with an identifier of the form \texttt{ТЗ-Ф-XX}, an automated check and the location of the corresponding test must be specified. Non-functional requirements are confirmed by separate test methods: \hseFill{methods for verifying non-functional requirements}. Additionally, the completeness of the supporting test materials attached to this document is verified.
\end{hseExample}

\clearpage
% Source pages are printed without the bottom stamp (as in real
% document sets) — only the sheet number and the document code at the top.
\pagestyle{appendix}

% ============================================================
% LIST OF REFERENCES
% ============================================================
\section*{LIST OF REFERENCES}
\addcontentsline{toc}{section}{LIST OF REFERENCES}

% Hint: add sources on the technologies of your project BEFORE the GOST
% standards, in alphabetical order, following the example:
% \bibitem{pytest}
% pytest Documentation. Accessed: Mar.~14, 2026. [Online]. Available: \url{https://docs.pytest.org/}
\renewcommand{\refname}{}
\begin{thebibliography}{99}
\bibitem{gost19101}
\emph{GOST 19.101-77. Types of Programs and Program Documents}, Unified System for Program Documentation. Moscow, Russia: IPK Izdatelstvo Standartov, 2001.

\bibitem{gost19103}
\emph{GOST 19.103-77. Designations of Programs and Program Documents}, Unified System for Program Documentation. Moscow, Russia: IPK Izdatelstvo Standartov, 2001.

\bibitem{gost19104}
\emph{GOST 19.104-78. Basic Inscriptions}, Unified System for Program Documentation. Moscow, Russia: IPK Izdatelstvo Standartov, 2001.

\bibitem{gost19105}
\emph{GOST 19.105-78. General Requirements for Program Documents}, Unified System for Program Documentation. Moscow, Russia: IPK Izdatelstvo Standartov, 2001.

\bibitem{gost19106}
\emph{GOST 19.106-78. Requirements for Program Documents Executed in Printed Form}, Unified System for Program Documentation. Moscow, Russia: IPK Izdatelstvo Standartov, 2001.

\bibitem{gost19301}
\emph{GOST 19.301-79. Test Program and Methodology. Requirements for the Content and Presentation}, Unified System for Program Documentation. Moscow, Russia: IPK Izdatelstvo Standartov, 2001.

\bibitem{gost19604}
\emph{GOST 19.604-78. Rules for Making Changes to Program Documents Executed in Printed Form}, Unified System for Program Documentation. Moscow, Russia: IPK Izdatelstvo Standartov, 2001.
\end{thebibliography}



\clearpage

% ============================================================
% ЛИСТ РЕГИСТРАЦИИ ИЗМЕНЕНИЙ
% ============================================================
\input{../common/change_log}

\end{document}
